\documentclass[reqno,11pt]{amsart}
\usepackage{math327-lecture}
\usepackage{needspace}
\lectureinfo{12}{Oct. 2, 2026}{Groups of Order Six and the Counting Formula}

% Based on handwritten Lecture 9, Sep. 22, 2025, together with the
% material on groups of order six moved from lecture-11.tex.
\begin{document}
\makelecturetitle

\section{Review: the counting formula}

Let $H\leq G$. Last time we saw that $G$ is partitioned into the left
cosets of $H$, and that $G/H$ denotes the set of \emph{distinct} left
cosets. If $R\subseteq G$ contains exactly one representative of each
left coset, then
\[
  G=\bigsqcup_{r\in R}rH.
\]
Since each coset $rH$ is in bijection with $H$, counting both sides gives
the following.

\begin{theorem}[Counting formula]\label{thm:counting}
If $G$ is finite and $H\leq G$, then
\begin{equation}
  |G|=|G/H|\cdot|H|=[G:H]\,|H|.
\end{equation}
In particular (Lagrange's theorem), if $x\in G$ then
$\operatorname{ord}(x)=|\langle x\rangle|$ divides $|G|$.
\end{theorem}

\section{Classifying groups of order six}

Let $|G|=6$. Lagrange's theorem (Theorem~\ref{thm:counting}) shows that an element can have order
$1$, $2$, $3$, or $6$. We first establish that elements of orders $2$ and
$3$ must occur.

\begin{lemma}
Every finite group of even order contains an element of order $2$.
\end{lemma}
\begin{proof}
Pair each element $a$ with its inverse $a^{-1}$. The distinct sets
$\{a,a^{-1}\}$ partition $G$ into parts of size one or two. A part has
size one exactly when $a=a^{-1}$.
If the identity were the only such element, then $|G|$ would be one plus
an even number, hence odd. Thus there is some $a\neq e_G$ with
$a=a^{-1}$. It satisfies $a^2=e_G$ and has order $2$.
\end{proof}

Here is a similar exercise from the homework: 

\begin{question} Show that every group of order $35$ has an element of order $5$ or $7$.

Hints: If there are no elements of order $5$, show that every element has order $7$ or $1$ (the identity). Now if $x$ and $y$ are elements of order $7$, show that the subgroups $\langle x \rangle $ and $\langle y \rangle$ are either equal or contain just the identity element in common. Now count!  
 \end{question}

\begin{lemma}
If $a^2=e_G$ for every $a\in G$, then $G$ is abelian.
\end{lemma}
\begin{proof}
Every element is its own inverse, so for $a,b\in G$,
\[
  ba=b^{-1}a^{-1}=(ab)^{-1}=ab.
\]
\end{proof}

\begin{proposition}
A group of order $6$ contains an element of order $3$.
\end{proposition}
\begin{proof}
Suppose first that every nonidentity element has order $2$.
The preceding lemma makes $G$ abelian. Choose distinct nonidentity
elements $a,b$. Then
\[
  K=\{e_G,a,b,ab\}=\langle a,b\rangle
\]
is a subgroup: since $a,b$ commute and $a^2=b^2=e_G$, every product of
powers of $a,b$ reduces to one of these four elements. They are distinct:
$ab=e_G$ would imply $a=b$, while $ab=a$ or $ab=b$ would force $b=e_G$
or $a=e_G$, respectively. Thus $|K|=4$, contradicting Lagrange's theorem
because $4$ does not divide $6$.

There is therefore an element $a$ whose order is neither $1$ nor $2$.
Its order is $3$ or $6$. In the first case, take $x=a$.
In the second case, take $x=a^2$; then $x^3=e_G$ and $x\neq e_G$, so
$x$ has order $3$.
\end{proof}

Choose $x$ of order $3$ and $y$ of order $2$, and put
$H=\langle x\rangle=\{e_G,x,x^2\}$. Since the elements of $H$ have
orders $1$ or $3$, we have $y\notin H$. The coset criterion gives
$yH\neq H$. Both cosets have three elements, so
\begin{equation}
  G=H\sqcup yH=\{e_G,x,x^2,y,yx,yx^2\}.
\end{equation}
In particular, $x$ and $y$ generate $G$.
\subsection*{How do \texorpdfstring{$x$}{x} and \texorpdfstring{$y$}{y} multiply?}
The product $xy$ is an element of $G$, so it must be one of the six
elements listed above. Which one?

First, $xy\notin H$. Indeed, if $xy\in H$, then
$y=x^{-1}(xy)\in x^{-1}H=H$, since $x^{-1}\in H$. This contradicts
$y\notin H$. Hence $xy\in yH=\{y,yx,yx^2\}$. Moreover $xy\neq y$, since
cancelling $y$ would give $x=e_G$. We are left with two possibilities:
\begin{equation}\label{eq:two-cases}
  xy=yx\qquad\text{or}\qquad xy=yx^2.
\end{equation}
A priori, both could happen.

\begin{proposition}\label{prop:mult-determined}
Once we know which case of \eqref{eq:two-cases} holds, the multiplication
in $G$ is completely determined.
\end{proposition}
\begin{proof}[Sketch] Let's just explain the idea by way of example: suppose $xy=yx^2$ and $x^3=y^2=1$. Now given two elements of $G = \{ 1 , x , x^2, y, yx, yx^2 \}$, we claim that we can always move all the $y$'s to the \textit{left} using the above rules. For example, \be{} x^2 \cdot (y x) = x \cdot ( x \cdot y) \cdot x = x \cdot (y x^2) \cdot x = x \cdot y = y x^2 \ee and the latter is one of the elements we listed in $G$. 
%	
%	
%	
%	
%Write $xy=yx^k$ with $k\in\{1,2\}$. Then
%\[
%  x^ay=x^{a-1}(xy)=x^{a-1}yx^{k}=\cdots=yx^{ka}
%\]
%for every $a\geq0$, by moving $y$ past one $x$ at a time. Every element of
%$G$ has the form $y^bx^a$ with $b\in\{0,1\}$ and $a\in\{0,1,2\}$. Using the
%rule above $c$ times,
%\[
%  (y^bx^a)(y^cx^d)=y^{b+c}x^{k^ca+d},
%\]
%and the exponents can be reduced using $y^2=e_G$ and $x^3=e_G$. Thus every
%product of two elements of $G$ is determined. For instance, when $k=2$,
%\[
%  (yx)\cdot y=y(xy)=y\cdot yx^2=x^2,\qquad
%  y\cdot(yx)=y^2x=x.
%\]
\end{proof}

Are there actually groups in which each case occurs? Yes:

\begin{itemize}
  \item \emph{Case $xy=yx$, $x^3=e_G$, $y^2=e_G$.} Now $x$ and $y$
  commute, so $G$ is abelian and $(xy)^n=x^ny^n$. Hence
  \[
    (xy)^2=x^2,\quad (xy)^3=y,\quad (xy)^4=x,\quad (xy)^5=x^2y,\quad
    (xy)^6=e_G.
  \]
  These six powers are distinct, so $xy$ has order $6$ and
  $G=\langle xy\rangle$ is the cyclic group $C_6$ of order $6$.
  For example, $\Z/6\Z$ with $x=\overline2$, $y=\overline3$.
  \item \emph{Case $xy=yx^2$, $x^3=e_G$, $y^2=e_G$.}  Note that $x y = yx^2$ is equivalent to (taking inverses of both sides and remembering that $y^{-1}=y$ and $x^{-1}=x^2$) \be{}  y x ^2 = y^{-1} x^{-1} = x^{-2} y^{-1} = x y. \ee But we have seen these relations before: this is just $S_3$:
  take $x=(1\,2\,3)$ and $y=(1\,3)$.
\end{itemize}

\begin{theorem}
Every group of order $6$ is isomorphic to exactly one of $C_6$ and $S_3$.
\end{theorem}
\begin{proof}
By the discussion above, a group of order $6$ has elements $x,y$ of orders
$3,2$ with $G=\{e_G,x,x^2,y,yx,yx^2\}$, and satisfies one of the relations
\eqref{eq:two-cases}. By Proposition~\ref{prop:mult-determined}, its
multiplication table is then the same as that of $C_6$ or of $S_3$,
respectively; matching $y^bx^a$ with the corresponding element gives an
isomorphism. The two groups are not isomorphic, since $C_6$ is abelian and
$S_3$ is not.
\end{proof}

%\section{Index in a tower of subgroups}
%
%Recall the counting formula $|G|=[G:H]\,|H|$, where the index
%$[G:H]=|G/H|$ is the number of left cosets of $H$ in $G$. (When we write
%$[G:H]$ we usually mean that it is finite.)
%
%\begin{theorem}[Tower law]\label{thm:tower}
%Let $K\leq H\leq G$ be a tower of subgroups. Then
%\begin{equation}
%  [G:K]=[G:H]\,[H:K].
%\end{equation}
%\end{theorem}
%\begin{proof}
%Assume $[G:H]=m$ and $[H:K]=n$ are finite. Then
%\[
%  G=g_1H\sqcup g_2H\sqcup\cdots\sqcup g_mH\quad(g_i\in G),
%  \qquad
%  H=h_1K\sqcup\cdots\sqcup h_nK\quad(h_j\in H).
%\]
%Substituting the second decomposition into the first,
%\[
%  G=g_1(h_1K\cup\cdots\cup h_nK)\cup\cdots\cup
%    g_m(h_1K\cup\cdots\cup h_nK)
%   =\bigcup_{\substack{1\leq i\leq m\\1\leq j\leq n}}g_ih_jK.
%\]
%It remains to show that the $mn$ cosets $g_ih_jK$ are distinct, and hence
%(being cosets) disjoint. Suppose $g_ih_jK=g_{i'}h_{j'}K$. Since
%$h_jK\subseteq H$, we have $g_ih_jK\subseteq g_iH$ and similarly
%$g_{i'}h_{j'}K\subseteq g_{i'}H$. So $g_iH\cap g_{i'}H\neq\varnothing$,
%which forces $i=i'$. Cancelling $g_i$ then gives $h_jK=h_{j'}K$, so
%$j=j'$. Therefore $G$ is the disjoint union of exactly $mn$ left cosets
%of $K$, and $[G:K]=mn$.
%\end{proof}
%
%When $G$ is finite, the tower law also follows directly from the counting
%formula:
%\[
%  [G:K]=\frac{|G|}{|K|}=\frac{|G|}{|H|}\cdot\frac{|H|}{|K|}=[G:H]\,[H:K].
%\]
%The proof above has the advantage that it works for infinite groups, as
%long as the indices are finite.
%
%\begin{example}
%For $8\Z\leq2\Z\leq\Z$ we have $[\Z:2\Z]=2$ and $[2\Z:8\Z]=4$ (the cosets
%are $8\Z$, $2+8\Z$, $4+8\Z$, $6+8\Z$), so
%$[\Z:8\Z]=2\cdot4=8$, matching Lecture~11.
%In $S_3$, the tower $\{1\}\leq\{1,x,x^2\}\leq S_3$ gives
%$6=[S_3:\{1\}]=2\cdot3$.
%\end{example}

\end{document}
